\documentclass[a4paper,11pt]{ltjsarticle}

% ===== 基本設定 =====
% ===== 基本設定 =====
\usepackage{luatexja}
\usepackage{amsmath, amssymb}
\usepackage{graphicx}
\usepackage{booktabs}
\usepackage{siunitx}
\usepackage{url}
\usepackage{hyperref}
\usepackage{listings}
\usepackage{xcolor}

% ===== SI単位設定 =====
\sisetup{
  detect-all,
  separate-uncertainty = true
}

% ===== コード表示設定 =====
\lstset{
  basicstyle=\ttfamily\small,
  frame=single,
  breaklines=true,
  numbers=left,
  numberstyle=\tiny,
  keywordstyle=\color{blue},
  commentstyle=\color{gray},
  stringstyle=\color{red}
}


% ===== レポート情報 =====
\title{実験レポートタイトル}
\author{氏名：山田 太郎 \\ 学籍番号：XXXXXXXX}
\date{\today}

\begin{document}

\maketitle

\section{目的}

本実験の目的を簡潔に記述する。

\section{原理}

実験で用いる理論や式を記述する。

例えば、測定値 $x$ の平均値 $\bar{x}$ は次式で表される。

\begin{equation}
  \bar{x} = \frac{1}{n}\sum_{i=1}^{n} x_i
\end{equation}

\section{実験方法}

使用した装置、測定手順、実験条件を記述する。

\section{結果}

測定結果を表や図で示す。

\begin{table}[htbp]
  \centering
  \caption{測定結果の例}
  \label{tab:result}
  \begin{tabular}{ccc}
    \toprule
    試行 & 時間 / \si{\second} & 距離 / \si{\meter} \\
    \midrule
    1 & 1.23 & 0.45 \\
    2 & 1.25 & 0.46 \\
    3 & 1.22 & 0.44 \\
    \bottomrule
  \end{tabular}
\end{table}

\section{考察}

結果から分かること、誤差の原因、理論値との比較などを記述する。

\section{結論}

実験で得られた結論を簡潔にまとめる。

\section{参考文献}

\begin{thebibliography}{9}
  \bibitem{example}
  著者名，『書籍名』，出版社，出版年．
\end{thebibliography}

\end{document}
